\documentclass[../main.tex]{subfiles}

\begin{document}

\section{Task 3: Vertical Climb, Burn, and Coast}\label{sec:task3}

\subsection{Switching Function and Thrust-Arc Structure}

Task~3 introduces a coast arc between the powered burn and orbital injection. Unlike Tasks~1 and~2, where the engine burns throughout at the fixed magnitude $T = cQ$, the engine may now shut down: the thrust magnitude becomes a bounded control $T \in [0, T_{\max}]$ (with $T_{\max} = cQ$), alongside the steering angle $\varphi$, and the mass flow is $\dot{m} = -T/c$. Collecting in the Hamiltonian the terms proportional to $T$ gives
\begin{equation}
    \mathcal{H} = \underbrace{\lambda_x v_x + \lambda_y v_y - \lambda_{v_y}}_{\mathcal{H}_0\ (\text{independent of } T)}
    \; + \; T\left(\frac{\lambda_{v_x}\cos\varphi + \lambda_{v_y}\sin\varphi}{m} - \frac{\lambda_m}{c}\right),
    \label{eq:H_grouped}
\end{equation}
where the mass term $-\lambda_m\dot{m} = -\lambda_m T/c$ now sits inside the $T$-bracket because $T$ is no longer constant. Inserting the optimal steering of the bilinear-tangent law, which maximizes the bracket by aligning the thrust with the velocity costate so that $\lambda_{v_x}\cos\varphi + \lambda_{v_y}\sin\varphi = \|\boldsymbol{\lambda}_v\|$ (Eq.~\ref{eq:maxcond}), the Hamiltonian becomes linear in the thrust magnitude:
\begin{equation}
    \mathcal{H} = \mathcal{H}_0 + T\,S(t), \qquad
    S(t) = \frac{\|\boldsymbol{\lambda}_v(t)\|}{m(t)} - \frac{\lambda_m(t)}{c}.
    \label{eq:switching}
\end{equation}
The coefficient $S(t)$ of $T$ is the \textbf{switching function}~\cite[Sec.~3.9]{bryson1975}: the marginal change of the Hamiltonian per unit thrust, i.e.\ the propulsive benefit $\|\boldsymbol{\lambda}_v\|/m$ (alignment of thrust with the primer vector $\boldsymbol{\lambda}_v$~\cite[Lecture~6]{dclvnotes}) net of the propellant penalty $\lambda_m/c$. By Pontryagin's maximum principle the optimal $T$ maximizes $\mathcal{H}$ pointwise over $[0, T_{\max}]$; since $\mathcal{H}$ is linear in $T$, the maximizer lies at a bound---the \textbf{bang--bang} law
\begin{equation}
    T^\ast = \begin{cases} T_{\max}, & S > 0 \quad(\text{engine on}),\\[2pt] 0,        & S < 0 \quad(\text{coast}). \end{cases}
    \label{eq:bangbang}
\end{equation}
A coast arc is therefore an interval where $S < 0$; there $T = 0$, so $\dot{m} = 0$ and the vehicle flies ballistically. Because $\dot{\lambda}_m = 0$ whenever $T = 0$ and the transversality condition \eqref{eq:trans_m} fixes $\lambda_m(t_f) = 1$, the mass costate stays constant at $\lambda_m \equiv 1$ along the coast; by continuity $\lambda_m(t_c) = 1$ at cutoff, so there the switching function reduces to $S(t_c) = \|\boldsymbol{\lambda}_v(t_c)\|/m(t_c) - 1/c$.

\textbf{Non-singularity of the coast switch.}\ The bang--bang law assumes that $S(t)$ vanishes only at isolated instants. If instead $S \equiv 0$ on a finite interval, the thrust would be singular, left undetermined by the first-order condition. That does not happen here: along the powered arc $\boldsymbol{\lambda}_v$ is affine in time, $m$ is positive and linear, and $\lambda_m$ is monotone (Eq.~\ref{eq:lm}), so $S = \|\boldsymbol{\lambda}_v\|/m - \lambda_m/c$ is real-analytic in time and can vanish only at isolated points. The thrust profile is therefore bang--bang: a single switch from $T_{\max}$ to $0$ at engine cutoff, with no singular sub-arc.

The three-phase mission sequence is:
\begin{enumerate}
    \item \textbf{Vertical climb}: $\varphi = \pi/2$ until $y = y_1 = 10^{-4}$ (identical to Task~2),
    \item \textbf{Powered burn}: bilinear-tangent steering until engine cutoff at time $t_c$,
    \item \textbf{Coast arc}: ballistic flight from $t_c$ to the injection point.
\end{enumerate}

\subsection{Engine Cutoff Conditions}

The shooting unknowns remain $\boldsymbol{\zeta} = [\lambda_{v_x,0},\; \lambda_{v_y,0},\; \lambda_y,\; \lambda_{m,0},\; t_\mathrm{burn}]^T$ (with time measured from the start of Phase~2). At engine cutoff the following five conditions must be satisfied simultaneously:
\begin{equation}
    \mathbf{r}_3(\boldsymbol{\zeta}) =
    \begin{bmatrix}
        v_{x,c} - 1 \\[2pt]
        y_c + \tfrac{1}{2}v_{y,c}^2 - y_f \\[2pt]
        \lambda_m(t_c) - 1 \\[2pt]
        \lambda_{v_y}(t_c) - \lambda_y\, v_{y,c} \\[2pt]
        S(t_c)
    \end{bmatrix} = \mathbf{0}.
    \label{eq:residual3}
\end{equation}

Conditions~1 and~2 are the terminal constraints $v_x(t_f) = 1$ and $y(t_f) = y_f$ carried back to the cutoff through the ballistic coast. The coast leaves $v_x$ unchanged, so the orbital horizontal speed must already be reached at cutoff: $v_{x,c} = 1$ (condition~1). The vertical motion is a parabola under constant gravity ($\dot v_y = -1$), so the coast reaches apogee ($v_y = 0$) at $\tau^\ast = v_{y,c}$ with altitude $y_c + \tfrac{1}{2}v_{y,c}^2$ (Section~\ref{sec:coast_arc}); imposing $y(t_f) = y_f$ gives condition~2. The same result is vertical energy conservation along the coast: with $g = 1$, $\tfrac{1}{2}v_y^2 + y$ is constant, so $\tfrac{1}{2}v_{y,c}^2 + y_c = y_f$ --- the residual vertical kinetic energy at cutoff is exactly what lifts the vehicle from $y_c$ to the target altitude. Condition~3 is the standard mass-costate transversality \eqref{eq:trans_m}. Condition~4 is the free-final-time condition $\mathcal{H} = 0$ \eqref{eq:trans_tf}, read on the coast. With $T = 0$ and $\lambda_x \equiv 0$ the coast Hamiltonian reduces to $\mathcal{H} = \lambda_y v_y - \lambda_{v_y}$, which is constant and equal to zero along the optimal trajectory (the same $\mathcal{H} \equiv 0$ I impose at $t_0$ in Task~1); evaluating it at cutoff, where $v_y = v_{y,c}$, gives $\lambda_{v_y}(t_c) = \lambda_y v_{y,c}$ directly. Equivalently, at apogee ($v_y = 0$) it reads $\lambda_{v_y}(t_f) = 0$, which propagates back to the cutoff through the linear evolution $\lambda_{v_y}(t_f) = \lambda_{v_y}(t_c) - \lambda_y v_{y,c}$. This is the Hamiltonian (free-time) condition, not the transversality of the constrained terminal $v_y$, whose costate is left free. Condition~5 is the bang-bang switching condition: the engine cuts off exactly when $S = 0$.

\subsection{Coast Arc}\label{sec:coast_arc}

During the ballistic coast ($T = 0$) the trajectory is analytically integrable. Denoting the elapsed coast time by $\tau$ and the cutoff state by subscript $c$:
\begin{equation}
    x(t_c + \tau) = x_c + v_{x,c}\,\tau, \qquad
    y(t_c + \tau) = y_c + v_{y,c}\,\tau - \tfrac{1}{2}\tau^2.
\end{equation}
The coast ends at apogee $\tau^* = v_{y,c}$, where $v_y = 0$ and
\begin{equation}
    y_f = y_c + \tfrac{1}{2}v_{y,c}^2,
\end{equation}
in agreement with condition~2 of \eqref{eq:residual3}. The vehicle mass is constant throughout the coast: $m_f = m_c$.

\subsection{Results}

Figure~\ref{fig:task3} shows the three-phase trajectory and mass profile for $Q = 2$, $y_f = 0.04$. The vertical climb (red) is short; the burn arc (blue) terminates below the target altitude; the ballistic coast (green dashed) ascends to the injection point. The mass profile confirms the model: it decreases linearly during both powered phases and remains constant during the coast. The final mass in Task~3 exceeds that of Task~2: by cutting the engine early and relying on the ballistic coast, less propellant is consumed while still satisfying all terminal constraints.

\begin{figure}[H]
    \centering
    \begin{subfigure}[b]{0.68\textwidth}
        \centering
        \includegraphics[width=\textwidth]{task3_task_3_trajectory}
        \caption{Trajectory: vertical climb (red), burn (blue), coast (green dashed).}
        \label{fig:task3_traj}
    \end{subfigure}
    \\[2ex]
    \begin{subfigure}[b]{0.68\textwidth}
        \centering
        \includegraphics[width=\textwidth]{task3_task_3_mass}
        \caption{Mass profile: linear decrease during powered phases, flat during coast.}
        \label{fig:task3_mass}
    \end{subfigure}
    \caption{Task 3: trajectory and mass profile for $Q = 2$, $y_f = 0.04$.}
    \label{fig:task3}
\end{figure}

Figure~\ref{fig:task3_angles} shows the thrust and velocity angles throughout the three phases. The dashed vertical lines mark the end of the vertical climb and the engine cutoff. During the coast no thrust angle is defined; the velocity angle $\psi$ continues to evolve ballistically until it reaches zero at apogee.

\begin{figure}[H]
    \centering
    \includegraphics[width=0.78\textwidth]{task3_task_3_angles}
    \caption{Task 3: thrust angle $\varphi$ and velocity angle $\psi$ across all three phases. Vertical dashed lines mark phase transitions.}
    \label{fig:task3_angles}
\end{figure}

A clear difference from Task~2 shows up in the thrust angle just after the vertical climb: $\varphi \approx 47^\circ$ here, against $\approx 71^\circ$ in Task~2.\footnote{The thrust-angle, flight-path-angle and velocity values quoted in this paragraph ($\varphi$, $\psi$, $v_{y,c}$) are read from the converged solutions computed by \texttt{main\_task2.m} and \texttt{main\_task3.m}, rather than measured off the figures.} Because the coast covers the upper half of the climb (from $y_c \approx 0.020$ up to $y_f = 0.040$) with the engine off, the burn no longer has to gain all the altitude on its own and can point the thrust closer to the horizontal from the start, that is, a more aggressive gravity turn that spends more of the thrust on horizontal acceleration. The same logic appears at cutoff: in Task~2 the thrust angle becomes negative ($\varphi \approx -27^\circ$), with the engine pushing downward in the last instants to null $v_y$ at injection, while here it stays positive ($\varphi \approx +20^\circ$) and the engine simply shuts off while $v_y$ is still positive ($v_{y,c} = 0.199$), letting gravity null it during the coast. Cutting off near-horizontally with some vertical velocity left over, instead of burning propellant to brake vertically, is what raises the final mass from $m_f = 0.112$ (Task~2) to $0.133$.

The velocity angle $\psi$ shows the same effect on the trajectory side. In Task~2 it falls from $90^\circ$ to a plateau near $37^\circ$ and then decreases monotonically to $0^\circ$ at injection, with the engine still firing. In Task~3 it instead drops to a much lower minimum ($\approx 5^\circ$), then rises back to $\approx 11^\circ$ at cutoff, and only reaches $0^\circ$ at apogee, with the engine off. The deep minimum is the counterpart of the smaller thrust angle: with $\varphi \approx 47^\circ$ and the mass still high just after the climb, the vertical thrust component $\tfrac{T}{m}\sin\varphi$ is smaller than $g$, so $v_y$ drops and the path flattens sharply; as propellant burns off, $\tfrac{T}{m}$ grows until $\tfrac{T}{m}\sin\varphi > g$ again, $v_y$ climbs back and $\psi$ recovers, which is how the burn builds up the residual $v_{y,c} = 0.199$ handed to the coast. In Task~2 the larger $\varphi$ keeps $\tfrac{T}{m}\sin\varphi > g$ from the start, so $v_y$ never collapses, $\psi$ holds its plateau, and the engine itself drives $\psi$ to zero. The U-shaped $\psi$ in Task~3, compared with the monotone $\psi$ in Task~2, is the signature of the coast on the trajectory.

\end{document}
